\section{Projective sublines}
Let $q\geq3$, $e\geq2$, $Q=q^e$, and let $\mathcal S$ consist of
$\Omega=PG(1,Q)$ and the $PGL(2,Q)$ images of $PG(1,q)$.
This is $S(3,q+1,Q+1)$. To verify uniqueness, carry an ordered triple to
$(\infty,0,1)$. Any image subline containing that triple can be
parametrized by a projectivity whose three preimages lie in $PG(1,q)$;
compose with its unique $PGL(2,q)$ normalization. The resulting
projectivity fixes three points and is the identity. Thus the unique
block containing the normalized triple is the standard subline.
The number of blocks through a pair is
\[
r=\frac{Q-1}{q-1}.
\]
In particular the parameters satisfy the firm-construction hypotheses
of [G, Theorem S.1].

\textbf{Theorem I.Subline.}
The geometry $\Gamma(\mathcal S)$ is firm and satisfies
$(\mathrm{IP})_2$ for every displayed $q,e$. It is residually connected
if and only if $e=2$. Among these connected geometries it is
flag-transitive under its full type-preserving automorphism group if
and only if $q=4$.

\emph{Proof.}
Firmness and $(\mathrm{IP})_2$ follow from [G, Theorems S.1 and S.6].
The derivative at $\infty$ is $AG(e,q)$: deleting $\infty$ from a block
through it gives $a\mathbb F_q+b$, $a\ne0$. Conversely every subline
through $\infty$ admits a representing projectivity fixing $\infty$,
by composing with a member of $PGL(2,q)$; it is therefore of this
form. Projective transitivity gives the same derivative at every point.
The common residue theorem S.5, Theorem B.AffinePlane and Theorem
B.SpaceObstruction now give connectedness exactly when $e=2$.
This proves the residual-connectivity assertion without identifying it
with flag transitivity, as the printed Lemma 5 of [L] does.

For $e=2$, the full design group is $P\Gamma L(2,q^2)$.
Here is the precise outside input for the upper bound. The design has
its sharply three-transitive $PGL(2,q^2)$ action, and the proof of [K,
Theorem 3, p.73] bounds its full three-transitive group by the semilinear
projective group. Its degree $q^2+1$ is neither a power of two nor one
of the Mathieu degrees in that classification. Conversely all
semilinear projectivities preserve the block family. The reconstruction
theorem S.2 identifies this full design group with the full
type-preserving group of $\Gamma(\mathcal S)$.

Write $q=p^f$. The chamber formula S.3 gives
\[
|\operatorname{Ch}\Gamma|=(q^2+1)q^2(q^2-1)q
 =|PGL(2,q^2)|q,
\qquad
|\operatorname{Aut}\Gamma|=2f|PGL(2,q^2)|.
\]
Transitivity requires $q\mid2f$. If $p\geq3$, then $p^f>2f$ for all
$f\geq1$. If $p=2$, the assumptions give $f\geq2$, and $2^f\leq2f$
only for $f=2$. Thus $q=4$ is necessary.

For sufficiency take $Q=16$. Normalize a flag $(x,T,B_0)$, $x\in T$
and $|T|=2$, to $x=\infty$, $T=\{\infty,0\}$,
$B_0=PG(1,4)$. This is possible using $PGL(2,16)$ and then the
$PGL(2,4)$ stabilizer of that block. Alternative blocks through $T$
are $\{\infty\}\cup\alpha\mathbb F_4$, where
$\alpha\mathbb F_4^*$ ranges over the four nonidentity classes of
$\mathbb F_{16}^*/\mathbb F_4^*$, a cyclic group of order five.
Frobenius $z\mapsto z^2$ cycles these four classes. Once a class is
chosen, the scalars in $\mathbb F_4^*$ act transitively on the three
possible copied points in $\alpha\mathbb F_4^*$. Hence the stabilizer
of the normalized flag is transitive on all chamber extensions, proving
sufficiency. \hfill$\square$

For completeness the six rank-two residues are described explicitly.
Use types $0,1,2,3$ for ordinary point, marked pair, block, copied point.
Fixing types $01,12,13$ gives respectively
$K_{r-1,q-1}$, $K_{2,q-1}$ and $K_{2,r-1}$.
Fixing $23$, namely $B_0,z\notin B_0$, leaves the subdivision of the
complete graph on $B_0$: each marked pair is uniquely determined by a
two-subset of $B_0$ and the unique block containing it together with $z$.
Fixing $03$ gives $AG(e,q)$ with one point and the lines through it
removed. Fixing $02$ gives its points outside a line $L$ and lines
meeting $L$ in exactly one point. The latter identifications follow by
taking the ordinary point to $\infty$ in the derivative just described;
the marked object records the residual line and its unique intersection
with $L$, or the residual point and its unique line through the removed
point, respectively. These constructions give inverse object maps and
preserve all incidences. The final residue has exactly
$(q^{e-1}-1)/(q-1)$ connected components by B.SpaceObstruction.

The literal $S(3,4,10)$ example, with $q=3,e=2$, is a finite
counterexample to the printed restriction $q=4$ for connectedness.
Its type sizes are $10,180,30,10$, and it has $2160$ chambers.
The exact checker verifies all $1881$ residues of rank at least two.
This computation checks a consequence of the uniform proof; it is not
the proof for arbitrary $q$ or $e$.
